% Additive mathematical audit and proof, 2026-09-18.
% Frozen providers are not edited. Analytic single-norm proof is supplied
% by the companion equilibrium and GPA audit appendices in this intake.
\begingroup
\section{The attained upper fibre and the complete native-row profile}
\label{sec:ATA}

\subsection{Original objects and exact meaning of every cutoff}
Keep the original simple quartet, its complete nonreduced conductor,
fixed period and logarithm branch.  In particular retain
\[
\begin{gathered}
 k=4l+1\ge9,\quad j=k+4,\quad q=(j+1)^2,\quad q'=(j-7)^2,
 \quad \Delta=q-q'=16j-48,\\
 c=j/2,\quad c'=j/2-4,\quad
 v=\operatorname{ord}_0E_A,\quad L=q-v,\quad m=\Delta-v>0,
 \quad s\in\{1,j\},\\
 E_A(z)=\sum_{a,b=0}^8a_{ab}e^{\beta_{8,a,b}z},\quad
 \beta_{d,a,b}=d/2+(2a-d)\delta+i(2b-d)\gamma,\\
 \mu_h=E_A^{(h)}(0),\quad \mu_v\ne0,\quad
 \chi_j(S)=\prod_{a,b=0}^j(S-\beta_{j,a,b}),\quad
 \chi_-(S')=\chi_{j-8}(S').
\end{gathered}\tag{ATA1}
\]
The actual conductor has $0\le v\le80$ by NIB1; no value of $v$
is replaced by zero.  For either centre $d=c,c'$ use precisely
\[
 \|P\|_{s,d}^2=\int_{\mathbb R}|P(d+iy)|^2dm_s(y),\qquad
 dm_s(y)=\frac{(2\pi)^{s/2}2^{s/2}}{4\pi\Gamma(s/2)}
             |\Gamma(s/4+iy/2)|^2dy.                 \tag{ATA2}
\]
Write $D^{\mathrm{nat}}_n$ for NTG18's form at native polynomial
degree $n$.  If a target cutoff $N$ is used, write instead
$D_N=D^{\mathrm{nat}}_{N-v}$, as in FTR3.  Thus the same row has
native degree $L+r$ and target cutoff $q+r$.  These conventions give
$D^{\mathrm{nat}}_{L+r}=D_{q+r}$, including the unchanged centre $c'$.

The original relations of NTG2 are
\[
 (\mathcal T_AP)(S')=\sum_{a,b=0}^8a_{ab}P(S'+\beta_{8,a,b}),
 \quad g_r=\mathcal T_A(\chi_jS^r),\quad
 \ell_r=\mu_v\binom{q+r}{v},\quad u_r=g_r/\ell_r.       \tag{ATA3}
\]
They are monic of degree $L+r$ and belong to
$\chi_-\mathbb C[S']_{\le m+r}$.  All $\ell_r$ are nonzero.
The full lower-ideal minimum and the attained relation minimum are
\[
\begin{aligned}
 \nu^-_{m+r,s}
 &=\min_{P\text{ monic},\ \deg P=m+r}\|\chi_-P\|_{s,c'}^2,\\
 \eta_{r,s}
 &=\min_{z\in\mathbb C^r}
       \left\|u_r+\sum_{t=0}^{r-1}z_tu_t\right\|_{s,c'}^2,
 \qquad \gamma_{r,s}=\log(\eta_{r,s}/\nu^-_{m+r,s}).
\end{aligned}\tag{ATA4}
\]
For $r=0$ the second minimum has its unique empty coefficient vector.
Both minima exist uniquely: in a finite polynomial coefficient space
the relevant Gram is strictly positive, because a nonzero polynomial
cannot vanish on a real interval and the original density is positive.
Completing its positive quadratic coefficient square gives the minimum.
NTG8 and NIB5--9 give these same minima, with no change of fibre.
Since every member of the second affine fibre belongs to the first
monic lower-ideal fibre, $\eta_{r,s}\ge\nu^-_{m+r,s}>0$.

\subsection{An affine bijection before the conductor and its attained slack}
Define the original upper minimum by
\[
 \nu^+_{r,s}=\min_{P\text{ monic},\ \deg P=r}
                      \|\chi_jP\|_{s,c}^2,
 \qquad P^+_{r,s}\text{ its unique minimizer}.           \tag{ATA5}
\]
For $P(S)=S^r+\sum_{t<r}p_tS^t$, linearity gives literally
\[
 \frac{\mathcal T_A(\chi_jP)}{\ell_r}
       =u_r+\sum_{t<r}p_t\frac{\ell_t}{\ell_r}u_t,
 \qquad p_t=z_t\frac{\ell_r}{\ell_t}.                  \tag{ATA6}
\]
These coefficient maps are mutually inverse affine maps between
the complete upper monic fibre and the complete original relation
fibre in ATA4.  Consequently
\[
 \eta_{r,s}=|\ell_r|^{-2}
       \min_{P\text{ monic},\ \deg P=r}
                         \|\mathcal T_A(\chi_jP)\|_{s,c'}^2.
                                                               \tag{ATA7}
\]
Let $E_r=\operatorname{span}(u_0,\ldots,u_{r-1})$, with $E_0=0$,
and let $\Pi_r$ be its orthogonal projection in ATA2.  Put
\[
 w_r=\mathcal T_A(\chi_jP^+_{r,s})/\ell_r,
 \qquad \tau_{r,s}=\|w_r\|_{s,c'}^2.
\]
Equation ATA6 gives $w_r-u_r\in E_r$, hence
$(I-\Pi_r)w_r=(I-\Pi_r)u_r$.  For every $e\in E_r$,
\[
 \|w_r+e\|_{s,c'}^2
 =\|(I-\Pi_r)u_r\|_{s,c'}^2+\|\Pi_rw_r+e\|_{s,c'}^2.
\]
Taking its exact minimum and then $e=0$ proves
\[
 \eta_{r,s}=\|(I-\Pi_r)u_r\|_{s,c'}^2,
 \qquad \tau_{r,s}=\eta_{r,s}+\|\Pi_rw_r\|_{s,c'}^2.    \tag{ATA8}
\]
For example if $G_{<r}=(\langle u_a,u_b\rangle)_{a,b<r}$,
$h_a=\langle u_a,w_r\rangle$, then
$\Pi_rw_r=\sum_{a<r}u_a(G_{<r}^{-1}h)_a$.
Thus this projection uses the entire original Gram inverse.
This is the complex Hermitian quadratic minimum underlying the
classical Schur calculation of Boyd and Vandenberghe
\cite[Appendix A.5.5, p.~650]{human:boyd-vandenberghe}; the displayed
coefficient map and square completion give its full specialization.

\subsection{A polynomial translation bound for the actual conductor}
The unchanged NTG3 polynomials and squared norms are
\[
 p_0=1,\quad p_1=y,\quad
 p_{n+1}=yp_n-n(n-1+s/2)p_{n-1},\qquad
 h_{n,s}=(2\pi)^{s/2}n!(s/2)_n.                          \tag{ATA9}
\]
Their generating function is
\[
 F_s(y,z)=\sum_{n\ge0}\frac{p_{n,s}(y)}{n!}z^n
       =(1+z^2)^{-s/4}e^{y\arctan z}.                  \tag{ATA10}
\]
Indeed multiplication of the recurrence by $z^n/n!$ and summation
gives $(1+z^2)\partial_zF_s=(y-(s/2)z)F_s$ and $F_s(y,0)=1$;
the right side of ATA10 solves that equation.  Its Taylor coefficients
therefore agree recursively, and the right side is analytic for
$|z|<1$.  The human-authored NIST DLMF generating formula
\cite[eq.~(18.23.7)]{human:dlmf18} gives the same identity under
$p_{n,s}(y)=n!P_n^{(s/4)}(y/2;\pi/2)$.  The previously verified
DLMF dictionary also supplies exactly ATA9's norm, including
$dy=2\,dx$ and $(2\pi)^{s/2}$; see NTG3 and LRP3.

For a complex shift $z_0$ expand
$e^{z_0\arctan z}=\sum_{a\ge0}e_a(z_0)z^a$.
Coefficient comparison gives
\[
 p_{n,s}(y+z_0)=\sum_{a=0}^n\frac{n!}{a!}
                       e_{n-a}(z_0)p_{a,s}(y).         \tag{ATA11}
\]
Fix an integer $d\ge1$ and $0\le a\le d$.  On $|z|=d/(d+1)$,
the power series of $\arctan$ gives
$|\arctan z|\le\operatorname{artanh}(d/(d+1))
=\tfrac12\log(2d+1)$.  Cauchy's coefficient estimate and
$(1+1/d)^d\le e$ give
\[
 |e_a(z_0)|\le e(2d+1)^{|z_0|/2}.                      \tag{ATA12}
\]
In the original orthonormal coordinates the remaining entry factor is
\[
 \frac{n!}{a!}\sqrt{\frac{h_{a,s}}{h_{n,s}}}
  =\sqrt{\prod_{b=a+1}^n\frac{b}{b-1+s/2}}
  \le\sqrt{2d+1}\qquad(0\le a\le n\le d, s\ge1).     \tag{ATA13}
\]
For the inequality, $s\ge1$ bounds the product by
$\prod_{b=a+1}^nb/(b-1/2)$, and each omitted factor $b/(b-1/2)$
exceeds one.  Its enlargement from $b=1$ to $n$ equals
$4^n/\binom{2n}{n}\le2n+1$: the middle binomial coefficient
is the maximum among the $2n+1$ coefficients with sum $4^n$.
At $n=0$ the product equals one.
The translation matrix has at most $(d+1)^2$ entries.  Its operator
norm is bounded by its Frobenius norm, so ATA12--13 prove, for every
polynomial of degree at most $d$,
\[
 \|p(\,\cdot+z_0)\|_s
 \le e(d+1)(2d+1)^{(|z_0|+1)/2}\|p\|_s.                \tag{ATA14}
\]

The centres in ATA1 give the exact shift
\[
 z_{ab}=-i(\beta_{8,a,b}-4)
       =(2b-8)\gamma-i(2a-8)\delta,
 \quad c'+iy+\beta_{8,a,b}=c+i(y+z_{ab}).
\]
Thus the triangle inequality in the unchanged measure, applied to
each literal conductor summand, gives
\[
\begin{gathered}
 \mathfrak C_A(d)=e(d+1)\sqrt{2d+1}
                \sum_{a,b=0}^8|a_{ab}|(2d+1)^{|z_{ab}|/2},\\
 \|\mathcal T_AP\|_{s,c'}\le\mathfrak C_A(d)\|P\|_{s,c}
       \quad(\deg P\le d),\\
 \eta_{r,s}\le\tau_{r,s}
       \le\frac{\mathfrak C_A(q+r)^2}{|\ell_r|^2}\nu^+_{r,s}.
\end{gathered}\tag{ATA15}
\]
For the fixed actual conductor the finite sum has fixed coefficients
and exponents.  It has a nonzero term since $\mu_v\ne0$; hence
$|\log\mathfrak C_A(d)|=O_{\mathrm{actual}}(\log(d+1))$.
Likewise the exact finite product in $\binom{q+r}{v}$ gives
$|\log|\ell_r||=O_{\mathrm{actual}}(\log q)$ uniformly on
$0\le r\le q+32j$.  No inverse-conductor bound is used.

\subsection{Receiving the complete finite analytic endpoint}
Use the actual root-comparison domain
\[
 j\ge257,\qquad j\sqrt{\delta^2+\gamma^2}\le2^{-33}q,
 \qquad R_j<r\le q+32j,
 \quad R_j=\lfloor j^{3/2}/\log j\rfloor.                \tag{ATA16}
\]
The additive GPA reconstruction ATG3--28 proves the complete monic comparisons
with $Q_+=q+(s-1)/4$ and $Q_-=q'+(s-1)/4$:
\[
 \nu^+_{r,s}\le\beta_se^{b^+_{q,r,s}}\mathfrak n(Q_+,r),
 \quad
 \nu^-_{m+r,s}\ge\beta_se^{b^-_{q',m+r,s}}\mathfrak n(Q_-,m+r),
 \quad\beta_s=\frac{(2\pi)^{s/2}}{\sqrt2\Gamma(s/2)}.
\]
The complete power-exponential norm proof EQ43--54 in
Section~\ref{eqreview:powerexponential} gives its finite endpoints
$H^-(Q,b)$ and $H^+(Q,b)$.
Define the actual endpoint by the full expression
\[
\begin{split}
 U_{j,s,r}={}&2\log\mathfrak C_A(q+r)-2\log|\ell_r|
 +b^+_{q,r,s}-b^-_{q',m+r,s}\\
 &+H^+(Q_+,r)-H^-(Q_-,m+r).
\end{split}\tag{ATA17}
\]
Here every factor belongs to the same original source order $s$.
The mass $\beta_s$ cancels only in this matched ratio.  Substitution
in ATA15 proves the two finite inequalities
\[
 0\le\gamma_{r,s}\le U_{j,s,r},\qquad
 0\le\log(\tau_{r,s}/\eta_{r,s})\le U_{j,s,r}-\gamma_{r,s}.
                                                               \tag{ATA18}
\]
In particular this explicitly defined $U_{j,s,r}$ is nonnegative.
The exact centre calculation EQ45--54 and AER11--14 proves uniformly on ATA16
\[
 U_{j,s,r}=\Delta\mathfrak f(r/q)
                  +O_{\mathrm{actual}}(q/r+\log q),
 \qquad
 \mathfrak f(t)=\log\frac{1+\rho_t}{1-\rho_t},\quad
 \frac{\rho_tK(\sqrt{1-\rho_t^2})}
      {E(\sqrt{1-\rho_t^2})}=\frac1{1+t}.               \tag{ATA19}
\]
The same calculation proves $\mathfrak f>0$, decreasing and smooth
on $(0,\infty)$, $\mathfrak f(t)=-\tfrac12\log t+O(\sqrt t)$
at zero, and $\sup_{0<t\le4}t|\mathfrak f'(t)|<\infty$.
The exact integral calculation PCB1--8, with the original EIQ energy
$F=-E_{\min}$ and logarithmic moment $\mathcal L$, gives
\[
 2\int_0^1\mathfrak f(t)dt
 =4\log(4/\pi)-10-2F(2,\pi)+4\mathcal L(2,\pi)
 =C_\partial.                                        \tag{ATA20}
\]
The right side is precisely the previously defined ECL16 coefficient.

\subsection{A finite discrepancy and weighted convergence of the actual rows}
For the original four-endpoint return retain
$\vartheta_0=\vartheta_q=1$ and $\vartheta_r=2$ for $0<r<q$.
HJR21--26, with its degree parameter set to $j$, supply the exact
total and its finite lower endpoint
\[
\begin{gathered}
 \mathcal T_{j,s}=\sum_{r=0}^q\vartheta_r\gamma_{r,s},\\
 \mathcal T^-_{j,s}
 =\max\{0,\mathcal F_{\partial,j}^{(s)}-e^-_{j,s}
       -\mathcal E^{\mathrm{FEX}}_{j,s}-E^{\mathrm{tr}}_{j,s}
       -(v-t_0)\mathcal L_s\}\le\mathcal T_{j,s},\\
 \mathcal T^-_{j,s}=\Delta qC_\partial+o(jq),\qquad
 \mathcal T_{j,s}=\Delta qC_\partial+o(jq).
\end{gathered}\tag{ATA21}
\]
Every error and original source restriction in this displayed endpoint
is the HJR25 expression, not a new independently chosen error.
The lower asymptotic follows from its displayed finite errors and the
same outer-return asymptotic used in HJR26; ATA20 gives $C_\partial>0$,
so its maximum with zero does not change that asymptotic.
NIB14 and NIB18--20 give the explicit prefix endpoint
\[
 E^{\mathrm{pre}}_{j,s}
 =\sum_{r=0}^{R_j}\vartheta_rE^{\mathrm{NIB}}_{j,s,r}
 =O_{\mathrm{actual}}(j^{5/2}+j^3/\log j)=o(jq).         \tag{ATA22}
\]
Define the finite number
\[
 \mathfrak D_{j,s}
 =\sum_{r=R_j+1}^q\vartheta_rU_{j,s,r}
                       -\mathcal T^-_{j,s}+E^{\mathrm{pre}}_{j,s}.
                                                               \tag{ATA23}
\]
All tail sums in this section end at $q$, even though ATA19 is proved
on the larger range for the separate displacement application.
Subtracting the actual prefix from the exact total gives
\[
\begin{split}
 0&\le\sum_{r=R_j+1}^q\vartheta_r(U_{j,s,r}-\gamma_{r,s})\\
  &=\sum_{r=R_j+1}^q\vartheta_rU_{j,s,r}
    -\mathcal T_{j,s}
    +\sum_{r=0}^{R_j}\vartheta_r\gamma_{r,s}
 \le\mathfrak D_{j,s}.
\end{split}\tag{ATA24}
\]
To evaluate its size, fix one constant for the uniform remainder
in ATA19.  Its summed absolute value is at most
\[
 C\sum_{r=R_j+1}^q(q/r+\log q)=O_{\mathrm{actual}}(q\log q)=o(jq).
\]
For a decreasing integrable positive function,
\[
 \int_{1/q}^{1+1/q}\mathfrak f(t)dt
 \le\frac1q\sum_{r=1}^q\mathfrak f(r/q)
 \le\int_0^1\mathfrak f(t)dt.
\]
Both bounds have the same limit.  The removed first $R_j$ terms have
sum at most $q\int_0^{R_j/q}\mathfrak f(t)dt=o(q)$, and the
single endpoint correction from its weight one is $O(1)$.
Therefore ATA19--20 imply
\[
 \sum_{r=R_j+1}^q\vartheta_rU_{j,s,r}
                      =\Delta qC_\partial+o(jq).
\]
Together with ATA21--22 this proves $\mathfrak D_{j,s}=o(jq)$.
For each row set $a_r=\Delta\mathfrak f(r/q)$; since $U_{j,s,r}\ge
\gamma_{r,s}$, the triangle inequality gives
$|\gamma_{r,s}-a_r|\le U_{j,s,r}-\gamma_{r,s}+|U_{j,s,r}-a_r|$.
Consequently
\[
\begin{gathered}
 \sum_{r=R_j+1}^q\vartheta_r
       |\gamma_{r,s}-\Delta\mathfrak f(r/q)|=o(jq),\\
 0\le\sum_{r=R_j+1}^q\vartheta_r
              \log(\tau_{r,s}/\eta_{r,s})
       \le\mathfrak D_{j,s}=o(jq).
\end{gathered}\tag{ATA25}
\]
The second line follows directly from ATA18 and ATA24; it preserves
the nonzero projection term in ATA8.  This convergence asserts a
weighted total of absolute differences.  It assigns no pointwise
lower estimate to a prescribed row.

\subsection{Correlated determinant and every sublinear row set}
NTG12--15 gives the compatible original columns $x_r$ and
$\Omega_r=I_m+\sum_{t=0}^rx_tx_t^*$, $\Omega_{-1}=I_m$.
For integers $-1\le R'<Q'\le q$ and
$Y=[x_{R'+1},\ldots,x_{Q'}]$, elimination of the entire earlier
relation Gram, as in NTG23--25, gives exactly
\[
 \log\det(I+Y^*\Omega_{R'}^{-1}Y)
 =\log\frac{\det\Omega_{Q'}}{\det\Omega_{R'}}
 =\sum_{r=R'+1}^{Q'}\gamma_{r,s}.                       \tag{ATA26}
\]
For completeness its graph-coordinate block Gram is
$I+Y^*Y-Y^*X_{R'}(I+X_{R'}^*X_{R'})^{-1}X_{R'}^*Y$;
multiplication verifies
$I-X_{R'}(I+X_{R'}^*X_{R'})^{-1}X_{R'}^*=\Omega_{R'}^{-1}$.
The leading later triangular block $H_{22}$ remains on both sides
of this Gram and cancels only against
$\prod_{r=R'+1}^{Q'}\nu^-_{m+r,s}$ in its determinant.
The earlier/later block $H_{12}$ disappears by the exact earlier
minimum, while all pairings in $Y^*\Omega_{R'}^{-1}Y$ remain.

ATA22 and ATA25 imply uniform convergence of cumulative sums:
\[
 \sup_{0\le n\le q}
 \left|\frac1{jq}\sum_{r=0}^n\gamma_{r,s}
             -16\int_0^{n/q}\mathfrak f(t)dt\right|\longrightarrow0.
                                                               \tag{ATA27}
\]
Here the true prefix has the upper bound ATA22.  Its comparison
profile contributes at most
$\Delta q\int_0^{R_j/q}\mathfrak f=o(jq)$.
The remaining total absolute error is bounded by ATA25 because
$\vartheta_r\ge1$.  Uniform Riemann-sum convergence follows from
the same decreasing-function bounds just proved, applied at every
upper endpoint.  Finally $\Delta/j\to16$.
For $R'/q\to a$, $Q'/q\to b$, $0\le a\le b\le1$, subtraction
of the two cumulative limits in ATA27 and ATA26 therefore proves
\[
 \frac{1}{jq}\log\det(I+Y^*\Omega_{R'}^{-1}Y)
                   \longrightarrow16\int_a^b\mathfrak f(t)dt.
                                                               \tag{ATA28}
\]
Adding the cases $Q'=q-1,q$ gives
\[
 \frac1{jq}\log\frac{\det\Omega_{q-1}\det\Omega_q}
                       {(\det\Omega_{R'})^2}
                   \longrightarrow32\int_a^1\mathfrak f(t)dt.
                                                               \tag{ATA29}
\]

Define a finite budget for each original set
$\mathscr I\subset\{0,\ldots,q\}$ by
\[
 \mathcal E_s(\mathscr I)
 =\sum_{\substack{r\in\mathscr I\\r\le R_j}}
       \vartheta_rE^{\mathrm{NIB}}_{j,s,r}
  +\sum_{\substack{r\in\mathscr I\\r>R_j}}
       \vartheta_rU_{j,s,r}.                            \tag{ATA30}
\]
If $h=|\mathscr I|$, monotonicity of $\mathfrak f$ and the uniform
remainder yield the explicit upper estimate
\[
 0\le\mathcal E_s(\mathscr I)
 \le E^{\mathrm{pre}}_{j,s}
   +2\Delta q\int_0^{h/q}\mathfrak f(t)dt
   +O_{\mathrm{actual}}(q\log q).                      \tag{ATA31}
\]
Indeed the largest sum of $h$ values of a decreasing sequence is its
first $h$ terms; $\sum_{r=1}^h\mathfrak f(r/q)$ is at most the
displayed integral times $q$.  The absolute remainder on any subset
is at most its whole-range sum.  Thus whenever $h=o(q)$,
$\mathcal E_s(\mathscr I)=o(jq)$, uniformly over the locations of
those $h$ rows, by absolute continuity of the integral at zero.
For $R'_j=\lfloor q/\log j\rfloor$, the expansion at zero gives
\[
 \mathcal E_s(\{0,\ldots,R'_j\})
 =O_{\mathrm{actual}}\left(j^{5/2}
          +\frac{j^3\log\log j}{\log j}+q\log q\right)=o(jq).
                                                               \tag{ATA32}
\]
In particular the $j^3/\log j$ term of ATA22 has been absorbed in
$j^3\log\log j/\log j$ for sufficiently large $j$, not discarded.

At every row the exact NTG22 update satisfies
$\Omega_r=\Omega_{r-1}+x_rx_r^*$ and
$x_r^*\Omega_{r-1}^{-1}x_r=e^{\gamma_{r,s}}-1$.
Cauchy--Schwarz in the $\Omega_{r-1}$ metric gives
$x_rx_r^*\preceq(e^{\gamma_{r,s}}-1)\Omega_{r-1}$.
After the unchanged conjugation and $A$ map of NTG20, it follows that
\[
 D_{q+r-1}\preceq D_{q+r}
        \preceq e^{\gamma_{r,s}}D_{q+r-1}
        \preceq e^{U_{j,s,r}}D_{q+r-1}.                 \tag{ATA33}
\]
For any original restriction or quotient pair, extend its actual
kernel and value frames to the whole original coefficient space.
Schur factorization decomposes the full determinant ratio into the
restriction, the specified quotient, and its exterior quotient.
Every factor is at least one, because an increasing metric increases
both restriction and a minimum over the identical affine value fibre.
The four-copy ambient logarithmic ratio is exactly $4\gamma_{r,s}$.
Therefore each original CPQ quotient increment is between zero and
$4\gamma_{r,s}$, with no quotient-rank multiplier.  The original
two chains in NIB16 give the signed row bound $4\gamma_{r,s}$ as
well.  Summing those same labelled occurrences over $\mathscr I$
proves
\[
\begin{gathered}
 0\le S_K^{\mathscr I},S_R^{\mathscr I}
       \le\mathcal T^{\mathscr I}\le\mathcal E_s(\mathscr I),\\
 |\mathcal L_e^{\mathscr I}|,|\mathcal P_e^{\mathscr I}|
          \le4\mathcal E_s(\mathscr I),\qquad
 0\le V_X^{\mathscr I}\le4\mathcal E_s(\mathscr I)
                       \quad(X=b,\mathrm{rem},s,A),\\
 4S_K^{\mathscr I}=\mathcal L_e^{\mathscr I}
       +V_b^{\mathscr I}+V_{\mathrm{rem}}^{\mathscr I}
       +V_s^{\mathscr I}+V_A^{\mathscr I},\qquad
 S_K^{\mathscr I}+S_R^{\mathscr I}=\mathcal T^{\mathscr I}.
\end{gathered}\tag{ATA34}
\]
All quantities in ATA34 are negligible at scale $jq$ on every
$o(q)$-row set.  This assertion concerns the identical row sums;
it does not commute, independently optimize, or identify the eight
programme flags.  Native-to-transported metric errors and both
threshold comparison errors remain those in the exact receiving
formulas FTR/IPR/STF.

\subsection{The independent proper source and all four displaced cutoffs}
Apply the preceding result separately at degree $k$, with its own
source orders $s_0=1,k$, centre $c'_k=k/2-4$, root count
$Q=q_k=(k+1)^2$ and root-comparison domain.  Put
\[
 b=q_{k+4}-q_k=8k+24,\qquad\Delta_k=16k-48.
\]
The proper form in PRD7--10 is
\[
 G_N=J_0^*D_{k,N-v}^{(s_0)}J_0,\qquad
 Q_N=\operatorname{Sch}(G_N;Z_E,Z_R),\qquad g(N)=\log\det Q_N.
                                                               \tag{ATA35}
\]
Its original frames $J_0,Z_E,Z_R$ are fixed.  Its four target-displaced
cutoffs in PRD22 are $Q+b-1,Q+b,2Q+2b-1,2Q+2b$ in the $N$
convention, and its own standard cutoffs are $Q-1,Q,2Q-1,2Q$.
Their exact difference of returns is consequently
\[
\begin{split}
 \delta V_Q={}&g(2Q+2b-1)-g(2Q-1)+g(2Q+2b)-g(2Q)\\
              &-g(Q+b-1)+g(Q-1)-g(Q+b)+g(Q).
\end{split}\tag{ATA36}
\]
Schur factorization in the single degree-$k$ ambient source, with
the actual restriction $J_0$ and actual quotient $Z_E$, proves
\[
 0\le g(Q+r)-g(Q+r-1)\le\gamma_{k,s_0,r}.              \tag{ATA37}
\]
This is the same frame-extension argument used before ATA34, applied
in one copy.  For a zero-dimensional value space its determinant is
one and the assertion also holds.
Expanding ATA36 by these increments gives its two negative intervals
$r=0,\ldots,b-1$ and $r=1,\ldots,b$, and its two positive intervals
$r=Q,\ldots,Q+2b-1$ and $r=Q+1,\ldots,Q+2b$.
Since $2b=16k+48\le32k$ for $k\ge3$, both positive intervals lie in
the extended range proved in ATA16--19 at degree $k$.
Their lower endpoints exceed $R_k$, so ATA18 applies there.
NIB14 applies at every low row, irrespective of its location relative
to $R_k$.  Therefore the following finite directed bounds hold:
\[
\begin{gathered}
 -E_Q^-\le\delta V_Q\le E_Q^+,\\
 E_Q^- =\sum_{r=0}^{b-1}E^{\mathrm{NIB}}_{k,s_0,r}
                   +\sum_{r=1}^{b}E^{\mathrm{NIB}}_{k,s_0,r},\\
 E_Q^+ =\sum_{r=Q}^{Q+2b-1}U_{k,s_0,r}
                   +\sum_{r=Q+1}^{Q+2b}U_{k,s_0,r}.
\end{gathered}\tag{ATA38}
\]
The low intervals contain $2b$ counted rows, each bounded by
$O_{\mathrm{actual}}(k\log k)$ by NIB15 for $r\le b=O(k)$.
Thus $E_Q^-=O_{\mathrm{actual}}(k^2\log k)$.
The two high intervals contain $4b$ counted rows.  On them
$r/Q=1+O(b/Q)$ and the derivative of $\mathfrak f$ is bounded
on a fixed neighbourhood of one.  Therefore ATA19 gives
\[
\begin{split}
 E_Q^+
 &=4b\Delta_k\mathfrak f(1)
       +O_{\mathrm{actual}}(\Delta_kb^2/Q+b\log Q)\\
 &=4b(16k-48)\mathfrak f(1)+O_{\mathrm{actual}}(k\log Q)
  =O_{\mathrm{actual}}(k^2).
\end{split}\tag{ATA39}
\]
No extra four-copy coefficient occurs in ATA37--39.  The coefficient
$4b$ counts the actual two high intervals of length $2b$.
It follows that $\delta V_Q=o(kq_k)$.  This estimates precisely the
displaced-cutoff correction and leaves the proper source's own
standard-window return in its exact degree-$k$ form.
In particular the separate target-prefix identity remains
\[
 (V_A^{>R}-V_Q)-(V_A-V_Q)=-V_A^{\le R}.                 \tag{ATA40}
\]
Neither this identity nor the profile determines the directional
allocation split, the signed eight-flag coefficient, or the transferred
arithmetic sign of the original projected action.
\endgroup
